\input{preamble}

\section*{Bhabha scattering 1}

Bhabha scattering is the process $e^-+e^+\rightarrow e^-+e^+$.

\begin{center}
\begin{tikzpicture}
\draw[dashed] (0,0) circle (0.5cm);
\draw[thick,->] (2,0) node[anchor=west] {$e^+$} -- (0.6,0);
\draw[thick,->] (-2,0) node[anchor=east] {$e^-$} -- (-0.6,0);
\draw[thick,->] (0.40,0.40) -- (1.3,1.3) node[anchor=south west] {$e^-$};
\draw[thick,->] (-0.4,-0.4) -- (-1.3,-1.3) node[anchor=north east] {$e^+$};
\draw (1,0.5) node {$\theta$};
\end{tikzpicture}
\end{center}

Momentum vectors where $E=\sqrt{p^2+m^2}$.
\begin{equation*}
p_1=\underset{e^- \, \longrightarrow}
{\begin{pmatrix}E\\0\\0\\p\end{pmatrix}}
\qquad
p_2=\underset{\longleftarrow \, e^+}
{\begin{pmatrix}E\\0\\0\\-p\end{pmatrix}}
\qquad
p_3=\underset{\substack{\phantom{e^-} \, \nearrow\\e^- \, \phantom{\nearrow}}}
{\begin{pmatrix}E\\p\sin\theta\cos\phi\\p\sin\theta\sin\phi\\p\cos\theta\end{pmatrix}}
\qquad
p_4=\underset{\substack{\phantom{\swarrow} \, e^+\\\swarrow \, \phantom{e^+}}}
{\begin{pmatrix}E\\-p\sin\theta\cos\phi\\-p\sin\theta\sin\phi\\-p\cos\theta\end{pmatrix}}
\end{equation*}

Spinors for $p_1$.
\begin{equation*}
u_{11}=\frac{1}{\sqrt{E+m}}
\underset{\text{spin up}}
{\begin{pmatrix}E+m\\0\\-p\\0\end{pmatrix}}
\qquad
u_{12}=\frac{1}{\sqrt{E+m}}
\underset{\text{spin down}}
{\begin{pmatrix}0\\E+m\\0\\p\end{pmatrix}}
\end{equation*}

Spinors for $p_2$.
\begin{equation*}
v_{21}=\frac{1}{\sqrt{E+m}}
\underset{\text{spin up}}
{\begin{pmatrix}p\\0\\E+m\\0\end{pmatrix}}
\qquad
v_{22}=\frac{1}{\sqrt{E+m}}
\underset{\text{spin down}}
{\begin{pmatrix}0\\-p\\0\\E+m\end{pmatrix}}
\end{equation*}

Spinors for $p_3$.
\begin{equation*}
u_{31}=\frac{1}{\sqrt{E+m}}
\underset{\text{spin up}}
{\begin{pmatrix}E+m\\0\\p_{3z}\\p_{3x}+ip_{3y}\end{pmatrix}}
\qquad
u_{32}=\frac{1}{\sqrt{E+m}}
\underset{\text{spin down}}
{\begin{pmatrix}0\\E+m\\p_{3x}-ip_{3y}\\-p_{3z}\end{pmatrix}}
\end{equation*}

Spinors for $p_4$.
\begin{equation*}
v_{41}=\frac{1}{\sqrt{E+m}}
\underset{\text{spin up}}
{\begin{pmatrix}p_{4z}\\p_{4x}+ip_{4y}\\E+m\\0\end{pmatrix}}
\qquad
v_{42}=\frac{1}{\sqrt{E+m}}
\underset{\text{spin down}}
{\begin{pmatrix}p_{4x}-ip_{4y}\\-p_{4z}\\0\\E+m\end{pmatrix}}
\end{equation*}

The scattering amplitude $\mathcal M$ is
\begin{equation*}
\mathcal M=\mathcal M_1-\mathcal M_2
\end{equation*}

where $\mathcal M_1$ is the amplitude for the $1\rightarrow3$, $4\rightarrow2$ process
\begin{equation*}
\mathcal M_1=\frac{\alpha}{t}
(\bar v_{2b}\gamma^\mu v_{4d})(\bar u_{3c}\gamma_\mu u_{1a})
\end{equation*}

and $\mathcal M_2$ is the amplitude for the $1\rightarrow2$, $4\rightarrow3$ process
\begin{equation*}
\mathcal M_2=\frac{\alpha}{s}
(\bar u_{3c}\gamma^\mu v_{4d})(\bar v_{2b}\gamma_\mu u_{1a})
\end{equation*}

The expected density $\langle|\mathcal M|^2\rangle$ is the average over spin states.
\begin{equation*}
\langle|\mathcal M|^2\rangle=\frac{1}{4}\sum_{abcd}|\mathcal M|^2
\end{equation*}

Expand the summand.
\begin{equation*}
\langle|\mathcal{M}|^2\rangle=\frac{1}{4}
\sum_{abcd}
\left(
\mathcal M_1\mathcal M_1^*-
\mathcal M_1\mathcal M_2^*-
\mathcal M_2\mathcal M_1^*+
\mathcal M_2\mathcal M_2^*
\right)
\end{equation*}

The summation results are
\begin{align*}
\sum_{abcd}
\mathcal M_1\mathcal M_1^*&=\frac{\alpha^2}{t^2}
\left(8 s^2 + 8 u^2 - 64 m^2 s - 64 m^2 u + 192 m^4\right)
\\
\sum_{abcd}
\mathcal M_1\mathcal M_2^*&=\frac{\alpha^2}{st}
\left(-8 u^2 + 64 m^2 u - 96 m^4\right)
\\
\sum_{abcd}
\mathcal M_2\mathcal M_2^*&=\frac{\alpha^2}{s^2}
\left(8 t^2 + 8 u^2 - 64 m^2 t - 64 m^2 u + 192 m^4\right)
\end{align*}

where $s$, $t$, and $u$ are the Mandelstam variables
\begin{align*}
s&=(p_1+p_2)^2
\\
t&=(p_1-p_3)^2
\\
u&=(p_1-p_4)^2
\end{align*}

For clarity define
\begin{align*}
f_{11}&=8 s^2 + 8 u^2 - 64 m^2 s - 64 m^2 u + 192 m^4
\\
f_{12}&=-8 u^2 + 64 m^2 u - 96 m^4
\\
f_{22}&=8 t^2 + 8 u^2 - 64 m^2 t - 64 m^2 u + 192 m^4
\end{align*}

Now write $\langle|\mathcal{M}|^2\rangle$ more compactly as
\begin{equation*}
\langle|\mathcal{M}|^2\rangle=\frac{\alpha^2}{4}
\left(\frac{f_{11}}{t^2}-\frac{2f_{12}}{st}+\frac{f_{22}}{s^2}\right)
\end{equation*}

For $E\gg m$ a useful approximation is to set $m=0$ and obtain
\begin{align*}
f_{11}&=8s^2+8u^2
\\
f_{12}&=-8u^2
\\
f_{22}&=8t^2+8u^2
\end{align*}

Then $\langle|\mathcal{M}|^2\rangle$ simplifies as
\begin{equation*}
\langle|\mathcal{M}|^2\rangle=2\alpha^2
\left(\frac{s^2+u^2}{t^2}+\frac{2u^2}{st}+\frac{t^2+u^2}{s^2}\right)
\end{equation*}

For $m=0$ the Mandelstam variables are
\begin{align*}
s&=4E^2
\\
t&=-2E^2(1-\cos\theta)
\\
u&=-2E^2(1+\cos\theta)
\end{align*}

Hence
\begin{equation*}
\langle|\mathcal{M}|^2\rangle
=2\alpha^2\left(
\frac{1+\cos^4(\theta/2)}{\sin^4(\theta/2)}
-\frac{2\cos^4(\theta/2)}{\sin^2(\theta/2)}
+\frac{1+\cos^2\theta}{2}
\right)
\end{equation*}

This can be written more compactly as
\begin{equation*}
\langle|\mathcal{M}|^2\rangle=\alpha^2\left(\frac{\cos^2\theta+3}{\cos\theta-1}\right)^2
\end{equation*}

The differential cross section is
\begin{equation*}
\frac{d\sigma}{d\Omega}=\frac{\langle|\mathcal{M}|^2\rangle}{4s}
=\frac{\alpha^2}{16E^2}\left(\frac{\cos^2\theta+3}{\cos\theta-1}\right)^2
\end{equation*}

Multiply by $(\hbar c)^2$ to convert from natural units to physical units.
\begin{equation*}
\frac{d\sigma}{d\Omega}=\frac{\alpha^2}{16E^2}
\left(\frac{\cos^2\theta+3}{\cos\theta-1}\right)^2
(\hbar c)^2
\end{equation*}

\href{https://georgeweigt.github.io/examples/bhabha-scattering-1-demo.html}{Eigenmath script}

\end{document}
